\documentclass[11pt]{article}
\newcommand{\ListaTitulo}{Lista 09 --- DBCA}
% Preâmbulo compartilhado — MATD48, Listas de Exercícios 2026
% Usado por Lista{NN}.tex e Gabarito{NN}.tex via % Preâmbulo compartilhado — MATD48, Listas de Exercícios 2026
% Usado por Lista{NN}.tex e Gabarito{NN}.tex via % Preâmbulo compartilhado — MATD48, Listas de Exercícios 2026
% Usado por Lista{NN}.tex e Gabarito{NN}.tex via \input{preamble}

\usepackage[utf8]{inputenc}
\usepackage[T1]{fontenc}
\usepackage[brazil]{babel}
\usepackage{fullpage}
\usepackage{amsmath,amssymb}
\usepackage{enumitem}
\usepackage{xcolor}
\usepackage{fancyhdr}
\usepackage{titlesec}
\usepackage{listings}
\usepackage{hyperref}
\usepackage{booktabs}
\usepackage{graphicx}

\definecolor{matd48azul}{HTML}{0B4F6C}
\definecolor{matd48verde}{HTML}{2E8B57}

\titleformat{\section}{\normalfont\Large\bfseries\color{matd48azul}}{\thesection}{1em}{}
\titleformat{\subsection}{\normalfont\large\bfseries\color{matd48azul}}{\thesubsection}{1em}{}

\providecommand{\ListaTitulo}{Lista de Exercícios}

\setlength{\headheight}{14pt}
\pagestyle{fancy}
\fancyhf{}
\lhead{\small MATD48 --- Planejamento de Experimentos A}
\rhead{\small \ListaTitulo}
\cfoot{\thepage}
\renewcommand{\headrulewidth}{0.4pt}

\lstset{
  basicstyle=\ttfamily\small,
  breaklines=true,
  frame=single,
  columns=fullflexible,
  backgroundcolor=\color{gray!8},
  commentstyle=\color{matd48verde},
  keywordstyle=\color{matd48azul}\bfseries,
  language=R,
  extendedchars=true,
  literate=%
    {á}{{\'a}}1 {à}{{\`a}}1 {â}{{\^a}}1 {ã}{{\~a}}1
    {é}{{\'e}}1 {ê}{{\^e}}1 {í}{{\'i}}1 {ó}{{\'o}}1
    {ô}{{\^o}}1 {õ}{{\~o}}1 {ú}{{\'u}}1 {ç}{{\c c}}1
    {Á}{{\'A}}1 {À}{{\`A}}1 {Â}{{\^A}}1 {Ã}{{\~A}}1
    {É}{{\'E}}1 {Ê}{{\^E}}1 {Í}{{\'I}}1 {Ó}{{\'O}}1
    {Ô}{{\^O}}1 {Õ}{{\~O}}1 {Ú}{{\'U}}1 {Ç}{{\c C}}1
}

\hypersetup{colorlinks=true, linkcolor=matd48azul, urlcolor=matd48azul, citecolor=matd48azul}

\newcommand{\cabecalho}[2]{%
  \begin{center}
    {\Large\bfseries\color{matd48azul} #1}\\[0.3em]
    {\large #2}\\[0.3em]
    \rule{\linewidth}{0.6pt}
  \end{center}
  \vspace{0.5em}
}

\newenvironment{questao}[1]{\vspace{1em}\noindent\textbf{Questão #1.}\hspace{0.4em}}{\par}
\newenvironment{solucao}{\vspace{0.3em}\noindent\textit{Solução.}\hspace{0.4em}\color{black!85}}{\par\vspace{0.5em}}


\usepackage[utf8]{inputenc}
\usepackage[T1]{fontenc}
\usepackage[brazil]{babel}
\usepackage{fullpage}
\usepackage{amsmath,amssymb}
\usepackage{enumitem}
\usepackage{xcolor}
\usepackage{fancyhdr}
\usepackage{titlesec}
\usepackage{listings}
\usepackage{hyperref}
\usepackage{booktabs}
\usepackage{graphicx}

\definecolor{matd48azul}{HTML}{0B4F6C}
\definecolor{matd48verde}{HTML}{2E8B57}

\titleformat{\section}{\normalfont\Large\bfseries\color{matd48azul}}{\thesection}{1em}{}
\titleformat{\subsection}{\normalfont\large\bfseries\color{matd48azul}}{\thesubsection}{1em}{}

\providecommand{\ListaTitulo}{Lista de Exercícios}

\setlength{\headheight}{14pt}
\pagestyle{fancy}
\fancyhf{}
\lhead{\small MATD48 --- Planejamento de Experimentos A}
\rhead{\small \ListaTitulo}
\cfoot{\thepage}
\renewcommand{\headrulewidth}{0.4pt}

\lstset{
  basicstyle=\ttfamily\small,
  breaklines=true,
  frame=single,
  columns=fullflexible,
  backgroundcolor=\color{gray!8},
  commentstyle=\color{matd48verde},
  keywordstyle=\color{matd48azul}\bfseries,
  language=R,
  extendedchars=true,
  literate=%
    {á}{{\'a}}1 {à}{{\`a}}1 {â}{{\^a}}1 {ã}{{\~a}}1
    {é}{{\'e}}1 {ê}{{\^e}}1 {í}{{\'i}}1 {ó}{{\'o}}1
    {ô}{{\^o}}1 {õ}{{\~o}}1 {ú}{{\'u}}1 {ç}{{\c c}}1
    {Á}{{\'A}}1 {À}{{\`A}}1 {Â}{{\^A}}1 {Ã}{{\~A}}1
    {É}{{\'E}}1 {Ê}{{\^E}}1 {Í}{{\'I}}1 {Ó}{{\'O}}1
    {Ô}{{\^O}}1 {Õ}{{\~O}}1 {Ú}{{\'U}}1 {Ç}{{\c C}}1
}

\hypersetup{colorlinks=true, linkcolor=matd48azul, urlcolor=matd48azul, citecolor=matd48azul}

\newcommand{\cabecalho}[2]{%
  \begin{center}
    {\Large\bfseries\color{matd48azul} #1}\\[0.3em]
    {\large #2}\\[0.3em]
    \rule{\linewidth}{0.6pt}
  \end{center}
  \vspace{0.5em}
}

\newenvironment{questao}[1]{\vspace{1em}\noindent\textbf{Questão #1.}\hspace{0.4em}}{\par}
\newenvironment{solucao}{\vspace{0.3em}\noindent\textit{Solução.}\hspace{0.4em}\color{black!85}}{\par\vspace{0.5em}}


\usepackage[utf8]{inputenc}
\usepackage[T1]{fontenc}
\usepackage[brazil]{babel}
\usepackage{fullpage}
\usepackage{amsmath,amssymb}
\usepackage{enumitem}
\usepackage{xcolor}
\usepackage{fancyhdr}
\usepackage{titlesec}
\usepackage{listings}
\usepackage{hyperref}
\usepackage{booktabs}
\usepackage{graphicx}

\definecolor{matd48azul}{HTML}{0B4F6C}
\definecolor{matd48verde}{HTML}{2E8B57}

\titleformat{\section}{\normalfont\Large\bfseries\color{matd48azul}}{\thesection}{1em}{}
\titleformat{\subsection}{\normalfont\large\bfseries\color{matd48azul}}{\thesubsection}{1em}{}

\providecommand{\ListaTitulo}{Lista de Exercícios}

\setlength{\headheight}{14pt}
\pagestyle{fancy}
\fancyhf{}
\lhead{\small MATD48 --- Planejamento de Experimentos A}
\rhead{\small \ListaTitulo}
\cfoot{\thepage}
\renewcommand{\headrulewidth}{0.4pt}

\lstset{
  basicstyle=\ttfamily\small,
  breaklines=true,
  frame=single,
  columns=fullflexible,
  backgroundcolor=\color{gray!8},
  commentstyle=\color{matd48verde},
  keywordstyle=\color{matd48azul}\bfseries,
  language=R,
  extendedchars=true,
  literate=%
    {á}{{\'a}}1 {à}{{\`a}}1 {â}{{\^a}}1 {ã}{{\~a}}1
    {é}{{\'e}}1 {ê}{{\^e}}1 {í}{{\'i}}1 {ó}{{\'o}}1
    {ô}{{\^o}}1 {õ}{{\~o}}1 {ú}{{\'u}}1 {ç}{{\c c}}1
    {Á}{{\'A}}1 {À}{{\`A}}1 {Â}{{\^A}}1 {Ã}{{\~A}}1
    {É}{{\'E}}1 {Ê}{{\^E}}1 {Í}{{\'I}}1 {Ó}{{\'O}}1
    {Ô}{{\^O}}1 {Õ}{{\~O}}1 {Ú}{{\'U}}1 {Ç}{{\c C}}1
}

\hypersetup{colorlinks=true, linkcolor=matd48azul, urlcolor=matd48azul, citecolor=matd48azul}

\newcommand{\cabecalho}[2]{%
  \begin{center}
    {\Large\bfseries\color{matd48azul} #1}\\[0.3em]
    {\large #2}\\[0.3em]
    \rule{\linewidth}{0.6pt}
  \end{center}
  \vspace{0.5em}
}

\newenvironment{questao}[1]{\vspace{1em}\noindent\textbf{Questão #1.}\hspace{0.4em}}{\par}
\newenvironment{solucao}{\vspace{0.3em}\noindent\textit{Solução.}\hspace{0.4em}\color{black!85}}{\par\vspace{0.5em}}


\begin{document}

\cabecalho{Lista de Exercícios 09}{Delineamento em blocos completos aleatorizados (Capítulo 4 / Aula 09)}

\noindent \textit{Entrega: até a próxima aula. Justifique todas as respostas — nestas listas, a
justificativa vale mais que a resposta final. O gabarito completo está em \texttt{Gabarito09.pdf}.}

\begin{questao}{1} \textbf{(Agricultura --- tabela de ANOVA do DBCA)}

Um agrônomo compara $t=4$ variedades de milho em $b=5$ blocos (talhões com fertilidade
diferente), medindo produtividade (sc/ha). O ajuste do DBCA produziu $SQ_{Total}=850$,
$SQ_{Trat}=320$ e $SQ_{Bloco}=210$.

\begin{enumerate}[label=(\alph*)]
  \item Complete a tabela de ANOVA (graus de liberdade, $SQ_{Erro}$, quadrados médios).
  \item Calcule a estatística $F_0$ para o teste do efeito de variedade.
  \item Sabendo que $F_{0{,}05;\,3;\,12}=3{,}49$, há evidência de diferença entre variedades ao
    nível de $5\%$? Justifique com base no valor de $F_0$.
\end{enumerate}
\end{questao}

\begin{questao}{2} \textbf{(Dedução --- decomposição ortogonal da soma de quadrados)}

Considere um DBCA balanceado com $t$ tratamentos e $b$ blocos (uma observação $Y_{ij}$ por
combinação tratamento-bloco). Defina
$$
a_i = \bar Y_{i\cdot} - \bar Y_{\cdot\cdot}, \qquad
b_j = \bar Y_{\cdot j} - \bar Y_{\cdot\cdot}, \qquad
e_{ij} = Y_{ij} - \bar Y_{i\cdot} - \bar Y_{\cdot j} + \bar Y_{\cdot\cdot},
$$
de modo que $Y_{ij}-\bar Y_{\cdot\cdot} = a_i + b_j + e_{ij}$.

\begin{enumerate}[label=(\alph*)]
  \item Mostre que $\sum_i a_i = 0$ e $\sum_j b_j = 0$.
  \item Mostre que, para cada $i$ fixo, $\sum_j e_{ij}=0$, e que, para cada $j$ fixo,
    $\sum_i e_{ij}=0$.
  \item Use os itens (a) e (b) para mostrar que os três produtos cruzados que aparecem ao
    elevar $Y_{ij}-\bar Y_{\cdot\cdot}$ ao quadrado e somar sobre $i,j$ se anulam, concluindo que
    $$
    SQ_{Total} = SQ_{Trat} + SQ_{Bloco} + SQ_{Erro}
    $$
    exatamente, sem termo de covariância entre tratamento e bloco.
\end{enumerate}
\end{questao}

\begin{questao}{3} \textbf{(Psicologia --- blocagem como aleatorização restrita)}

Uma psicóloga quer comparar três técnicas de relaxamento (respiração guiada, relaxamento
muscular progressivo, mindfulness breve) quanto à redução de ansiedade. Ela classifica os 30
pacientes recrutados em três blocos por gravidade basal da ansiedade (leve, moderada, grave; 10
pacientes por bloco) e, \emph{dentro de cada bloco}, sorteia os pacientes entre as três técnicas.

\begin{enumerate}[label=(\alph*)]
  \item Na notação de resultados potenciais, escreva o efeito causal médio que compara a técnica
    de respiração guiada ($i=1$) ao mindfulness breve ($i=3$).
  \item Explique por que sortear \emph{dentro} de cada bloco de gravidade (em vez de sortear entre
    todos os 30 pacientes ignorando a gravidade) ainda produz um estimador não viesado desse
    efeito causal. Que papel exato a aleatorização (agora restrita a cada bloco) desempenha nessa
    garantia?
  \item Sob quais circunstâncias a blocagem por gravidade basal reduz o erro-padrão da comparação
    entre técnicas, em vez de apenas gastar graus de liberdade do erro sem retorno?
\end{enumerate}
\end{questao}

\begin{questao}{4} \textbf{(Agricultura --- dado faltante, fórmula de Yates)}

Em um DBCA com $t=5$ tratamentos e $b=4$ blocos, a produtividade de uma parcela foi perdida
(praga localizada). Excluindo essa observação: o total do tratamento correspondente é
$T'=58$, o total do bloco correspondente é $B'=78$, e o total geral é $G'=380$.

\begin{enumerate}[label=(\alph*)]
  \item Estime o valor perdido pela fórmula de Yates.
  \item Após substituir o valor estimado e recalcular a ANOVA, quantos graus de liberdade tem o
    erro (em relação ao DBCA completo, sem dado faltante)? Justifique.
\end{enumerate}
\end{questao}

\begin{questao}{5} \textbf{(Eficiência relativa)}

Em um DBCA com $t=6$ tratamentos e $b=5$ blocos, obteve-se $QM_{Bloco}=45$ e $QM_{Erro}=15$.

\begin{enumerate}[label=(\alph*)]
  \item Calcule $gl_{Bloco}$, $gl_{Erro}$ (do DBCA) e $gl_{Erro}$ que um DCA hipotético teria
    tido com o mesmo número de unidades.
  \item Calcule a eficiência relativa $ER(DBCA,DCA)$.
  \item Interprete o valor numérico: quantas vezes mais repetições um DCA precisaria para igualar
    a precisão deste DBCA? Compare qualitativamente com o exemplo do viveiro de pepineiros visto
    em aula, em que $ER\approx1{,}02$.
\end{enumerate}
\end{questao}

\begin{questao}{6} \textbf{(Dissertativa --- matriz de delineamento)}

A matriz de delineamento de um DBCA pode ser escrita como $X=[\mathbf 1\mid X_{trat}\mid
X_{bloco}]$, sem nenhuma coluna de produto entre $X_{trat}$ e $X_{bloco}$.

\begin{enumerate}[label=(\alph*)]
  \item O que, exatamente, essa ausência de colunas de produto representa em termos do modelo
    estatístico do DBCA?
  \item Por que essa mesma estrutura de $X$ explica por que não se deve usar o teste $F$ sobre o
    termo de bloco para decidir se ele deve permanecer no modelo?
\end{enumerate}
\end{questao}

\begin{questao}{7} \textbf{(Ciência de dados --- fatores confundidos vs. blocagem)}

Um analista quer avaliar se um novo algoritmo de recomendação (\texttt{Algoritmo}: novo vs. antigo)
aumenta o tempo médio de sessão em um site. Por decisão de produto, o algoritmo novo foi ativado
\emph{apenas} para usuários cadastrados depois de 1º de janeiro de 2026; usuários cadastrados antes
dessa data continuam recebendo o algoritmo antigo, sem exceção.

\begin{enumerate}[label=(\alph*)]
  \item Identifique o segundo fator que está \textbf{completamente confundido} com
    \texttt{Algoritmo} nesse desenho, e explique por quê.
  \item Explique por que \emph{nenhuma quantidade adicional} de dados coletados sob essa mesma
    regra resolveria o problema — em particular, o que aconteceria ao tentar ajustar
    \texttt{lm(tempo\_sessao \textasciitilde{} Algoritmo + <o fator do item (a)>)}.
  \item Reformule o desenho usando o princípio de blocagem para que o efeito do algoritmo possa
    ser estimado sem confusão, identificando explicitamente o que faria o papel de bloco e como o
    tratamento seria sorteado dentro de cada bloco.
\end{enumerate}
\end{questao}

\end{document}
